\chapter{The Fast Fourier Transform}

In the last chapter we built the Discrete Fourier Transform (DFT): the machine that
takes $N$ samples of a signal and tells you, for each of $N$ frequencies, how much of
that frequency is present and in what phase. Geometrically, the DFT projects your
signal vector onto $N$ rotating phasors and reads off the coordinates. It is the right
idea. The problem is that, done the obvious way, it is \emph{slow} -- so slow that for
a long time people simply did not compute spectra of real signals at all.

Then, in 1965, two researchers published a four-page note describing an algorithm that
computed the very same numbers, exactly, but vastly faster. It did not approximate the
DFT; it reorganized the arithmetic so that enormous amounts of redundant work
collapsed. Spectra that would have taken a day took a second. Whole fields -- radar,
seismology, medical imaging, digital audio, wireless -- were rebuilt on top of it. The
algorithm is the \emph{Fast Fourier Transform} (FFT), and it is one of the most
consequential pieces of arithmetic ever written down.

This chapter is about \emph{why} the FFT is fast and \emph{how} it works, built from
the geometry up. We will see exactly where the wasted work in the naive DFT hides, watch
divide-and-conquer expose it, draw the famous butterfly, understand the strange
bit-reversal shuffle, and then -- crucially -- learn to actually \emph{use} an FFT:
zero-padding, reading the output of a real-valued FFT, and what the numbers mean. By the
end you will both understand the FFT and be able to call one correctly, which matters
because our wrist gesture controller will run an FFT on a tiny microcontroller hundreds
of times per second.

\section{Recalling the DFT, and counting its cost}
\label{ch09:dft-cost}

Let $x[0], x[1], \dots, x[N-1]$ be our $N$ samples. The DFT produces $N$ complex
numbers $X[0], \dots, X[N-1]$ by
\begin{equation}
\label{ch09:dft-defn}
X[k] \;=\; \sum_{n=0}^{N-1} x[n]\, e^{-\jj\, 2\pi k n / N},
\qquad k = 0, 1, \dots, N-1 .
\end{equation}
Each $X[k]$ asks: ``how much does the signal align with the phasor that completes
exactly $k$ whole turns over the window?'' The factor $e^{-\jj 2\pi k n/N}$ is that
phasor, sampled at instant $n$.

That exponential appears so often that it earns a name. Define the \textbf{twiddle
factor}
\begin{equation}
\label{ch09:twiddle}
W_N \;=\; e^{-\jj\, 2\pi / N},
\qquad\text{so that}\qquad
X[k] \;=\; \sum_{n=0}^{N-1} x[n]\, W_N^{\,kn}.
\end{equation}
$W_N$ is a single small clockwise rotation: it is the $N$th root of unity sitting just
below the positive real axis on the unit circle. All the phasors in the DFT are just
powers of this one number. That is the secret the FFT will exploit, but first let us see
the cost of ignoring it.

\begin{intuitionbox}
The DFT is a matrix-vector product. Stack the outputs into a vector
$\vec{X}$ and the inputs into $\vec{x}$; then $\vec{X} = \mat{F}\,\vec{x}$, where the
$N\times N$ matrix $\mat{F}$ has entries $F_{kn} = W_N^{kn}$. Every row of $\mat{F}$ is
one phasor; multiplying by $\mat{F}$ measures the overlap of $\vec{x}$ with each. A
dense $N\times N$ matrix times a vector takes about $N^2$ multiply-add operations. That
single sentence is the entire performance problem.
\end{intuitionbox}

\begin{focusbox}
\textbf{The naive DFT is $O(N^2)$.} Each of the $N$ outputs $X[k]$ is a sum over $N$
inputs, so it needs $N$ complex multiplications and $N-1$ complex additions. With $N$
outputs that is $N^2$ complex multiplies in total. Double $N$ and the work
\emph{quadruples}.
\end{focusbox}

To feel why $O(N^2)$ hurts, put numbers on it. A complex multiply is a handful of real
multiplies and adds; call the whole DFT ``$N^2$ operations'' for scale.

\begin{center}
\begin{tabular}{r r r r}
\toprule
$N$ & $N^2$ (naive) & $N\log_2 N$ (FFT) & speedup \\
\midrule
$64$       & $4{,}096$        & $384$           & $\approx 11\times$ \\
$1{,}024$  & $1{,}048{,}576$  & $10{,}240$      & $\approx 102\times$ \\
$65{,}536$ & $\approx 4.3\times 10^{9}$ & $\approx 1.0\times 10^{6}$ & $\approx 4096\times$ \\
$10^{6}$   & $10^{12}$        & $\approx 2\times 10^{7}$ & $\approx 50000\times$ \\
\bottomrule
\end{tabular}
\end{center}

A one-million-point spectrum is $10^{12}$ operations the naive way -- minutes to hours on
a real machine -- versus $2\times10^{7}$ with the FFT, a small fraction of a second. The
gap is not a constant factor you can shrug off; it grows without bound. Figure
\ref{ch09:fig-growth} shows the two curves pulling apart.

\begin{figure}[H]
\centering
\begin{tikzpicture}
\begin{axis}[
  width=11.5cm, height=6.4cm,
  xlabel={transform size $N$}, ylabel={operations (relative)},
  domain=2:256, samples=60,
  ymode=log, xmode=log, log basis x=2,
  xmin=2, xmax=256, ymin=1, ymax=80000,
  grid=both, grid style={gray!18},
  legend pos=north west, legend cell align=left,
  axis lines=box,
]
\addplot[specred, very thick, samples=60] {x^2};
\addlegendentry{naive DFT $\sim N^2$}
\addplot[sigblue, very thick, samples=60] {x*log2(x)};
\addlegendentry{FFT $\sim N\log_2 N$}
\addplot[gray, dashed, thick, samples=60] {x};
\addlegendentry{linear $N$ (reference)}
\end{axis}
\end{tikzpicture}
\caption{Operation count versus $N$ on log--log axes. The naive DFT (red) climbs as
$N^2$; the FFT (blue) hugs the nearly linear $N\log_2 N$ curve. The vertical gap
\emph{between} them keeps widening, which is why the FFT is not a minor optimization but
a change in what is computable.}
\label{ch09:fig-growth}
\end{figure}

\begin{connectionbox}
This is exactly the divide-and-conquer story from the algorithms companion. Sorting went
from $O(N^2)$ (insertion/selection sort) to $O(N\log N)$ (mergesort) by splitting the
problem in half, solving the halves, and stitching the results together cheaply. The FFT
is the same recursion applied to the DFT: split, recurse, combine. If you have met the
Master Theorem, the FFT satisfies $T(N) = 2\,T(N/2) + O(N)$, whose solution is
$O(N\log N)$. The whole intellectual content of this chapter is finding the cheap
``combine'' step.
\end{connectionbox}

\section{Where the redundancy hides}
\label{ch09:redundancy}

Before we can remove wasted work we have to \emph{see} it. The naive DFT recomputes the
same phasor values over and over and never reuses partial sums between different outputs.
Two structural facts about $W_N$ are the key.

\begin{whybox}
\textbf{Two symmetries of the twiddle factor.} Because $W_N = e^{-\jj 2\pi/N}$ lives on
the unit circle:
\begin{itemize}
\item \emph{Periodicity:} $W_N^{\,k+N} = W_N^{\,k}$. Going one extra full turn returns
you to the same point. So there are really only $N$ distinct phasor values, even though
the index $kn$ in \eqref{ch09:twiddle} ranges up to $(N-1)^2$.
\item \emph{Half-turn (the crucial one):} $W_N^{\,k + N/2} = -W_N^{\,k}$, because half a
revolution is a sign flip: $e^{-\jj\pi} = -1$. And $W_N^{2} = W_{N/2}$: skipping every
other point on an $N$-circle lands you exactly on an $(N/2)$-circle.
\end{itemize}
These tiny identities are the entire engine of the FFT. The first says values repeat; the
second says a size-$N$ problem secretly \emph{contains} a size-$N/2$ problem. Figure
\ref{ch09:fig-twiddle} shows the half-turn symmetry on the unit circle.
\end{whybox}

\begin{figure}[H]
\centering
\begin{tikzpicture}[>=Stealth, scale=2.1, every node/.style={font=\footnotesize}]
\draw[gray!45] (-1.25,0) -- (1.25,0);
\draw[gray!45] (0,-1.25) -- (0,1.25);
\draw[gray!55,thick] (0,0) circle (1);
% the eight powers of W_8 (clockwise from +1)
\foreach \k in {0,...,7}{
  \pgfmathsetmacro{\ang}{-45*\k}
  \fill[sigblue] (\ang:1) circle (0.9pt);
}
% highlight W^1 and W^5 = -W^1 (half-turn apart)
\draw[->,very thick,quotientorange] (0,0) -- (-45:1);
\draw[->,very thick,specred] (0,0) -- (135:1);
\node[anchor=north west] at (-45:1.02) {$W_8^{1}$};
\node[anchor=south east] at (135:1.02) {$W_8^{5}=-W_8^{1}$};
\node[anchor=south west] at (0:1.02) {$W_8^{0}=1$};
\node[anchor=north] at (-90:1.05) {$W_8^{2}$};
\node[anchor=east] at (180:1.04) {$W_8^{4}=-1$};
\end{tikzpicture}
\caption{The eight twiddle factors $W_8^{k}=e^{-\jj 2\pi k/8}$ are equally spaced on the
unit circle, stepping clockwise. A power and its half-turn partner sit diametrically
opposite: $W_8^{5}=-W_8^{1}$ (orange and red). That single sign flip is the half-turn
identity $W_N^{k+N/2}=-W_N^{k}$ that lets one twiddle multiply produce two outputs in the
butterfly.}
\label{ch09:fig-twiddle}
\end{figure}

The naive computation throws all this away. When it computes $X[k]$ and then $X[k+N/2]$,
it grinds through two full independent sums even though, term by term, their phasors
differ only by a sign. Multiply that wasted symmetry across all $N$ outputs and you have
found the factor that the FFT recovers.

\section{Divide and conquer: split by even and odd}
\label{ch09:split}

Assume $N$ is a power of two (we will return to other sizes). The Cooley--Tukey idea is
to split the \emph{input} samples into two interleaved halves -- the even-indexed
samples and the odd-indexed samples -- compute a DFT of each half, and then combine. Let
\begin{equation}
\label{ch09:eo-defn}
e[m] = x[2m], \qquad o[m] = x[2m+1], \qquad m = 0,\dots, \tfrac{N}{2}-1 .
\end{equation}
Now split the sum in \eqref{ch09:twiddle} into its even and odd terms:
\begin{align}
X[k]
&= \sum_{n \text{ even}} x[n]\,W_N^{kn} \;+\; \sum_{n \text{ odd}} x[n]\,W_N^{kn} \notag\\
&= \sum_{m=0}^{N/2-1} x[2m]\,W_N^{k(2m)}
   \;+\; \sum_{m=0}^{N/2-1} x[2m+1]\,W_N^{k(2m+1)} .
\label{ch09:split-1}
\end{align}
Look at the exponents. In the even sum, $W_N^{k(2m)} = (W_N^{2})^{km} = W_{N/2}^{km}$ --
the size-$N$ twiddle became a size-$N/2$ twiddle, exactly by the identity $W_N^2 =
W_{N/2}$. In the odd sum, pull one factor of $W_N^{k}$ out front:
$W_N^{k(2m+1)} = W_N^{k}\,W_{N/2}^{km}$. Substituting,
\begin{equation}
\label{ch09:split-2}
X[k] \;=\; \underbrace{\sum_{m=0}^{N/2-1} e[m]\,W_{N/2}^{km}}_{\displaystyle E[k]}
\;+\; W_N^{k}\,\underbrace{\sum_{m=0}^{N/2-1} o[m]\,W_{N/2}^{km}}_{\displaystyle O[k]}.
\end{equation}

\begin{happeningbox}
Read \eqref{ch09:split-2} slowly: $E[k]$ is \emph{the DFT of the even samples} and
$O[k]$ is \emph{the DFT of the odd samples}, each a transform of length $N/2$. We have
rewritten one size-$N$ DFT in terms of two size-$N/2$ DFTs plus a single twiddle
multiply $W_N^k$ on the odd half. That is the divide-and-conquer split: a big projection
expressed as two small projections, glued by a rotation.
\end{happeningbox}

There is a catch and it is a gift. $E[k]$ and $O[k]$ are defined for $k = 0, \dots,
N/2-1$ only, because they are length-$N/2$ transforms. But $X[k]$ needs $k$ all the way
to $N-1$. How do we get the top half of the output from half-length transforms? The
half-turn identity. For the top half, write $k + N/2$ with $k$ in the lower range:
\begin{align}
X[k]        &= E[k] + W_N^{k}\,O[k], \label{ch09:butterfly-lo}\\
X[k + N/2]  &= E[k] - W_N^{k}\,O[k], \qquad k = 0, \dots, \tfrac{N}{2}-1 .
\label{ch09:butterfly-hi}
\end{align}
The minus sign in \eqref{ch09:butterfly-hi} is the half-turn $W_N^{k+N/2} = -W_N^k$
doing its work; and $E$, $O$ are periodic with period $N/2$, so $E[k+N/2]=E[k]$ and
likewise for $O$. \emph{Two outputs come out of one twiddle multiply.} The product
$W_N^k\,O[k]$ is computed once and then added and subtracted. Half the multiplications
just vanished -- and that is only the top level of the recursion.

\begin{walkbox}
Why does splitting the input ``by even and odd'' work so cleanly, while splitting it into
``first half / second half'' does not? Because the DFT couples each output to every
input through the \emph{product} $kn$ in the exponent. Striding by $2$ in $n$ turns
$W_N^{2}$ into $W_{N/2}$ -- a clean change of scale on the unit circle. Splitting by
contiguous blocks does not produce that scaling; it produces a phase shift, which leads
to a different (and also valid) factorization called decimation-in-frequency. Even/odd is
decimation-in-time, the version drawn most often.
\end{walkbox}

\section{The butterfly}
\label{ch09:butterfly}

Equations \eqref{ch09:butterfly-lo}--\eqref{ch09:butterfly-hi} describe a tiny two-input,
two-output computation that is the atom of every FFT. It is called a \emph{butterfly}
because of how the data-flow diagram looks. You feed it $E[k]$ and $O[k]$; you multiply
the odd one by the twiddle $W_N^k$; then you form a sum and a difference.

\begin{figure}[H]
\centering
\begin{tikzpicture}[>=Stealth, every node/.style={font=\small}]
% nodes
\node (e) at (0,1.4) [circle,draw,fill=fillsig!40,minimum size=6mm] {$E[k]$};
\node (o) at (0,-1.4)[circle,draw,fill=fillsig!40,minimum size=6mm] {$O[k]$};
\node (xa) at (6,1.4) [circle,draw,fill=sigblue!22,minimum size=6mm] {$X[k]$};
\node (xb) at (6,-1.4)[circle,draw,fill=sigblue!22,minimum size=6mm] {$X[k{+}\tfrac{N}{2}]$};
% twiddle multiply node on the odd line
\node (w) at (2.6,-1.4) [rectangle,draw,fill=quotientorange!25,minimum size=6mm] {$\times\,W_N^{k}$};
% sum/diff nodes
\node (plus)  at (4.4,1.4)  [circle,draw,fill=white,minimum size=5mm] {$+$};
\node (minus) at (4.4,-1.4) [circle,draw,fill=white,minimum size=5mm] {$-$};
% edges
\draw[->,thick,sigblue] (e) -- (plus);
\draw[->,thick,sigblue] (e) -- (minus);
\draw[->,thick,quotientorange] (o) -- (w);
\draw[->,thick,quotientorange] (w) -- (plus);
\draw[->,thick,quotientorange] (w) -- (minus);
\draw[->,thick,sigblue] (plus) -- (xa);
\draw[->,thick,sigblue] (minus) -- (xb);
\end{tikzpicture}
\caption{The radix-2 decimation-in-time butterfly. The odd input is rotated by the
twiddle $W_N^k$; the result is added to the even input to make the lower output and
subtracted to make the upper output. One complex multiply, two complex adds, two outputs.
The whole FFT is thousands of these wired together.}
\label{ch09:fig-butterfly}
\end{figure}

\begin{focusbox}
\textbf{Cost of one butterfly:} exactly one complex multiply ($W_N^k\,O[k]$) and two
complex adds. There are $N/2$ butterflies at each stage and $\log_2 N$ stages, giving
$\tfrac{N}{2}\log_2 N$ complex multiplies overall -- the $O(N\log N)$ promise made
concrete.
\end{focusbox}

\subsection*{Stacking the stages}

We do not stop after one split. Each half-length DFT is itself a DFT of a power-of-two
size, so we split it again, and again, until the transforms have length~$1$. A length-$1$
DFT is the identity -- a single sample is its own spectrum, since the sum in
\eqref{ch09:dft-defn} has one term and the phasor is $W_1^0 = 1$. So the recursion
bottoms out for free.

Counting the levels: from $N$ down to $1$ by repeated halving takes $\log_2 N$ steps.
Each level does $O(N)$ work in butterflies. That is the recurrence
$T(N) = 2T(N/2) + O(N)$ once more, total $O(N\log N)$. Figure \ref{ch09:fig-stages}
shows the data flow for $N=8$: three stages of butterflies turning eight inputs into
eight outputs.

\begin{figure}[H]
\centering
\begin{tikzpicture}[>=Stealth, x=1.55cm, y=0.78cm, every node/.style={font=\footnotesize}]
% column x positions
\def\xin{0} \def\xa{1.5} \def\xb{3.0} \def\xc{4.5} \def\xout{6.0}
% input labels (bit-reversed order) at left
\foreach \y/\lab in {7/0,6/4,5/2,4/6,3/1,2/5,1/3,0/7}{
  \node[anchor=east] at (\xin,\y) {$x[\lab]$};
  \fill[sigblue] (\xin+0.12,\y) circle (1.6pt);
}
% nodes at each stage (8 rows)
\foreach \x in {\xa,\xb,\xc,\xout}{
  \foreach \y in {0,...,7}{ \fill[gray!55] (\x,\y) circle (1.3pt); }
}
% Stage 1: pairs (0,1)(2,3)(4,5)(6,7) -- adjacent rows
\foreach \p in {0,2,4,6}{
  \pgfmathtruncatemacro{\q}{\p+1}
  \draw[sigblue] (\xin+0.12,\p) -- (\xa,\p);
  \draw[sigblue] (\xin+0.12,\q) -- (\xa,\q);
  \draw[quotientorange] (\xin+0.12,\p) -- (\xa,\q);
  \draw[quotientorange] (\xin+0.12,\q) -- (\xa,\p);
}
% Stage 2: pairs differ by 2: (0,2)(1,3)(4,6)(5,7)
\foreach \p in {0,1,4,5}{
  \pgfmathtruncatemacro{\q}{\p+2}
  \draw[sigblue] (\xa,\p) -- (\xb,\p);
  \draw[sigblue] (\xa,\q) -- (\xb,\q);
  \draw[quotientorange] (\xa,\p) -- (\xb,\q);
  \draw[quotientorange] (\xa,\q) -- (\xb,\p);
}
% Stage 3: pairs differ by 4: (0,4)(1,5)(2,6)(3,7)
\foreach \p in {0,1,2,3}{
  \pgfmathtruncatemacro{\q}{\p+4}
  \draw[sigblue] (\xb,\p) -- (\xc,\p);
  \draw[sigblue] (\xb,\q) -- (\xc,\q);
  \draw[quotientorange] (\xb,\p) -- (\xc,\q);
  \draw[quotientorange] (\xb,\q) -- (\xc,\p);
}
% to output
\foreach \y/\lab in {7/0,6/1,5/2,4/3,3/4,2/5,1/6,0/7}{
  \draw[gray!60] (\xc,\y) -- (\xout,\y);
  \node[anchor=west] at (\xout+0.05,\y) {$X[\lab]$};
}
% stage labels
\node[anchor=south] at (\xa,7.7) {stage 1};
\node[anchor=south] at (\xb,7.7) {stage 2};
\node[anchor=south] at (\xc,7.7) {stage 3};
\end{tikzpicture}
\caption{An 8-point radix-2 FFT, decimation-in-time. Inputs enter on the left in
\emph{bit-reversed} order ($0,4,2,6,1,5,3,7$); three stages of butterflies (orange =
cross term carrying a twiddle, blue = straight term) produce the outputs in natural
order on the right. Each vertical pair of crossing lines is one butterfly from Figure
\ref{ch09:fig-butterfly}.}
\label{ch09:fig-stages}
\end{figure}

\section{Bit-reversal: why the inputs come in scrambled}
\label{ch09:bitreversal}

You may have noticed the inputs in Figure \ref{ch09:fig-stages} are not in order; they
are $x[0], x[4], x[2], x[6], x[1], x[5], x[3], x[7]$. This is not a typo, and it is not
optional. It falls directly out of the even/odd splitting, and understanding it removes
all the mystery.

Track what the recursion does to the indices. The first split sends even indices
$\set{0,2,4,6}$ to one group and odd $\set{1,3,5,7}$ to the other. Within ``even'' we
split again by even/odd \emph{of the half-index}: $\set{0,4}$ and $\set{2,6}$. One more
split isolates each element. Follow any index through these choices and you are reading
its binary digits from the least significant bit upward:

\begin{center}
\begin{tabular}{c c c c}
\toprule
index & binary & bits reversed & new position \\
\midrule
$0$ & $000$ & $000$ & $0$ \\
$1$ & $001$ & $100$ & $4$ \\
$2$ & $010$ & $010$ & $2$ \\
$3$ & $011$ & $110$ & $6$ \\
$4$ & $100$ & $001$ & $1$ \\
$5$ & $101$ & $101$ & $5$ \\
$6$ & $110$ & $011$ & $3$ \\
$7$ & $111$ & $111$ & $7$ \\
\bottomrule
\end{tabular}
\end{center}

\begin{happeningbox}
\textbf{Bit-reversal in one sentence:} repeatedly sorting by ``even or odd'' is the same
as sorting by the bits of the index from least-significant to most-significant, and that
is the same as reversing the bits. Sample $x[n]$ must start in the slot whose index is
$n$ with its binary digits reversed. For $n=1=001_2$ the reverse is $100_2 = 4$, so
$x[1]$ goes to slot $4$ -- exactly what the diagram shows.
\end{happeningbox}

This gives the standard \emph{in-place} FFT a clean shape: first permute the array by
bit-reversal, then run $\log_2 N$ stages of butterflies that read and write the same
array. No extra storage for the data, and the outputs emerge in natural frequency order.
(Some libraries instead take natural-order input and produce bit-reversed output; the
permutation has to live somewhere.)

\begin{trapbox}
Bit-reversal only produces a clean reshuffle when $N$ is a power of two. For other sizes,
$N$ is factored differently (a length-$12$ transform might split as $3\times 4$), the
``digit-reversal'' is in a mixed radix, and the butterflies are larger. Good libraries
(FFTW, the FFT in NumPy/SciPy, KissFFT, CMSIS-DSP) handle all this; you almost never
write bit-reversal yourself. But knowing it exists explains why FFT speed and even
\emph{accuracy} can depend sharply on whether $N$ is $1024$ versus $1023$ or $1000$.
\end{trapbox}

\begin{historybox}
The 1965 paper by James \textbf{Cooley} (IBM) and John \textbf{Tukey} (Princeton/Bell
Labs), ``An algorithm for the machine calculation of complex Fourier series,'' set off the
explosion -- partly because it arrived just as digital computers became widespread, and
partly, in the Cold War context, because fast spectra mattered for detecting nuclear
tests from seismic data. But the core idea is far older. Carl Friedrich \textbf{Gauss}
wrote down an equivalent doubling trick around \textbf{1805} -- before Fourier's own 1807
work on the series that bear his name -- to interpolate the orbits of the asteroids Pallas
and Juno. Gauss's note was in neglected Latin and unpublished in his lifetime; the method
was reinvented in pieces by Runge, Danielson, Lanczos and others over the next 150 years.
The lesson the field took to heart: the FFT was ``the most important numerical algorithm
of our lifetime'' (Gilbert Strang), and it had been hiding in plain sight for a century
and a half. See Cooley \& Tukey, \emph{Math.\ Comp.}\ \textbf{19} (1965), and Heideman,
Johnson \& Burrus, ``Gauss and the history of the FFT'' (1984).
\end{historybox}

\section{It computes the same numbers -- a small proof and a check}
\label{ch09:exact}

The FFT is \emph{exact}: up to floating-point rounding it returns the identical $X[k]$ as
\eqref{ch09:dft-defn}. We already proved the inductive step -- \eqref{ch09:split-2} and
the butterfly \eqref{ch09:butterfly-lo}--\eqref{ch09:butterfly-hi} are algebraic
identities, not approximations -- and the base case (length-$1$ DFT is the identity) is
immediate. By induction on $\log_2 N$, the recursion computes the DFT for every
power-of-two $N$. Let us verify by hand on the smallest interesting case.

\begin{example}[A 4-point FFT by hand]
\label{ch09:ex-n4}
Take $x = (1,\,2,\,0,\,1)$, so $N=4$ and $W_4 = e^{-\jj\pi/2} = -\jj$. The powers we need
are $W_4^0=1$, $W_4^1=-\jj$.

\emph{Split.} Even samples $e=(x[0],x[2])=(1,0)$; odd samples $o=(x[1],x[3])=(2,1)$.

\emph{Two length-2 DFTs} (a length-2 DFT is just sum and difference, since $W_2 = -1$):
\[
E = (\,e[0]+e[1],\ e[0]-e[1]\,) = (1,\,1), \qquad
O = (\,o[0]+o[1],\ o[0]-o[1]\,) = (3,\,1).
\]
\emph{Combine with butterflies}, using $W_4^0=1$, $W_4^1=-\jj$, for $k=0,1$:
\begin{align*}
X[0] &= E[0] + W_4^0 O[0] = 1 + 3 = 4, &
X[2] &= E[0] - W_4^0 O[0] = 1 - 3 = -2,\\
X[1] &= E[1] + W_4^1 O[1] = 1 + (-\jj)(1) = 1 - \jj, &
X[3] &= E[1] - W_4^1 O[1] = 1 + \jj .
\end{align*}
So $X = (4,\ 1-\jj,\ -2,\ 1+\jj)$. Check against the direct sum: $X[0]=\sum x[n]=4$;
good. And $X[3]=\overline{X[1]}$ as it must for a real input (next section). Plugging into
\eqref{ch09:dft-defn} directly gives the same four numbers, with a fraction of the
arithmetic of the brute-force $4^2=16$ multiplies.
\end{example}

\section{Using a real FFT: reading the output}
\label{ch09:reading}

Knowing how the FFT works is half the battle; reading its output correctly is the other
half, and it trips up nearly everyone at first. Suppose you sample a real signal at rate
$f_s$ (samples per second) and take an $N$-point FFT. Here is what the $N$ output numbers
mean.

\begin{focusbox}
\textbf{Frequency of bin $k$.} Output $X[k]$ measures the component at frequency
\begin{equation}
\label{ch09:bin-freq}
f_k \;=\; k\,\frac{f_s}{N} \quad\text{Hz}, \qquad k = 0, 1, \dots, N-1 .
\end{equation}
The spacing between bins is $\Delta f = f_s/N$, the \emph{frequency resolution}. Bin $0$
is DC (the average / zero frequency); bin $1$ is one cycle per window; and so on.
\end{focusbox}

\begin{trapbox}
\textbf{The top half is negative frequencies.} For $k > N/2$, bin $k$ represents the
frequency $f_k - f_s$, i.e.\ a \emph{negative} frequency $-(N-k)f_s/N$. Equivalently,
$X[N-k] = \overline{X[k]}$ for a real input -- the spectrum is conjugate-symmetric, so
the upper half is a mirror image of the lower half carrying no new information. Bin
$N/2$ is the Nyquist frequency $f_s/2$. If you plot all $N$ magnitudes you will see a
mirror; only bins $0$ through $N/2$ are independent for real data.
\end{trapbox}

Because of that symmetry, real-input FFTs have a dedicated, cheaper routine
(\texttt{numpy.fft.rfft}) that returns only the $\floor{N/2} + 1$ non-redundant bins,
$k = 0 \dots N/2$. Use it for sensor data -- it is faster and removes the temptation to
double-count energy. Figure \ref{ch09:fig-spectrum} shows a worked magnitude spectrum.

\begin{figure}[H]
\centering
\begin{tikzpicture}
\begin{axis}[
  width=12cm, height=5.6cm,
  xlabel={frequency (Hz)}, ylabel={magnitude $\abs{X[k]}$ (a.u.)},
  xmin=0, xmax=64, ymin=0, ymax=1.15,
  ytick={0,0.5,1.0},
  xtick={0,8,12,24,32,50,64},
  grid=both, grid style={gray!16},
  axis lines=box,
  ybar, bar width=2pt,
]
% f_s = 128 Hz, N = 128, so bin spacing = 1 Hz; plot only 0..64 (rfft range).
\addplot+[draw=sigblue, fill=fillsig] coordinates {
 (0,0.30)(1,0.03)(2,0.02)(3,0.04)(4,0.02)(5,0.03)(6,0.02)(7,0.03)
 (8,0.04)(9,0.05)(10,0.06)(11,0.18)(12,1.00)(13,0.20)(14,0.05)(15,0.03)
 (16,0.02)(17,0.03)(18,0.02)(19,0.03)(20,0.02)(21,0.03)(22,0.04)(23,0.12)
 (24,0.50)(25,0.13)(26,0.04)(27,0.03)(28,0.02)(29,0.03)(30,0.02)(31,0.03)
 (32,0.02)(33,0.03)(34,0.02)(35,0.03)(36,0.02)(37,0.03)(38,0.02)(39,0.03)
 (40,0.02)(41,0.03)(42,0.02)(43,0.03)(44,0.02)(45,0.03)(46,0.02)(47,0.03)
 (48,0.02)(49,0.03)(50,0.07)(51,0.03)(52,0.02)(53,0.03)(54,0.02)(55,0.03)
 (56,0.02)(57,0.03)(58,0.02)(59,0.03)(60,0.02)(61,0.03)(62,0.02)(63,0.02)(64,0.02)
};
\end{axis}
\end{tikzpicture}
\caption{Magnitude spectrum (\texttt{rfft}) of a $1\,$s record sampled at $f_s=128$ Hz,
$N=128$, so each bin is $1$ Hz wide. Two real tones appear as peaks at $12$ Hz
(amplitude $1$) and $24$ Hz (amplitude $0.5$); the bar at $0$ Hz is a DC offset, and the
small clutter is broadband noise. Only bins $0$ to $N/2 = 64$ are shown because the rest
are the conjugate mirror.}
\label{ch09:fig-spectrum}
\end{figure}

\subsection*{Amplitude calibration}

A raw FFT bin is \emph{not} the amplitude of the sinusoid. For a real cosine of
amplitude $A$ that lands exactly on bin $k$ (with $0 < k < N/2$), the FFT puts magnitude
$A\,N/2$ in bin $k$ (and the same in its mirror $N-k$). To recover $A$ from a one-sided
spectrum, scale by $2/N$:
\begin{equation}
\label{ch09:amp-cal}
A \;\approx\; \frac{2}{N}\,\bigl|X[k]\bigr| \quad (0 < k < N/2),
\qquad
A_{\text{DC}} = \frac{1}{N}\,\abs{X[0]},
\qquad
A_{\text{Nyq}} = \frac{1}{N}\,\abs{X[N/2]} .
\end{equation}
The DC and Nyquist bins use $1/N$ rather than $2/N$ because they have no separate mirror
partner. If you forget the $2/N$ you will report amplitudes off by a factor of $N/2$ --
the single most common FFT-reading mistake.

\section{Zero-padding and the difference between resolution and interpolation}
\label{ch09:zeropad}

A frequent move is \emph{zero-padding}: appending zeros to your signal before the FFT,
usually to reach a power-of-two length or to get more output points. It is useful, but it
is widely misunderstood, so let us be precise about what it does and does not buy.

\begin{whybox}
Zero-padding $x$ from length $N$ to length $M > N$ does \textbf{not} add information --
the zeros carry none. What it does is evaluate the \emph{same} underlying
continuous-frequency spectrum (the DTFT of your finite record) at more, finely spaced
points. The true frequency resolution -- your ability to \emph{separate} two nearby
tones -- is fixed by the duration of the real (nonzero) data, $\Delta f \approx f_s/N$.
Zero-padding gives a smoother, denser plot (interpolation), not sharper peaks. Two tones
unresolved at length $N$ stay unresolved no matter how many zeros you add.
\end{whybox}

\begin{intuitionbox}
Think of the finite record as the true infinite spectrum convolved with a
$\sinc$-shaped window response whose width is set by the record length. Zero-padding
samples that fixed sinc shape more densely. You see the shape better -- handy for
locating a peak precisely or for making a pretty curve -- but you cannot make the sinc
itself narrower by padding. To narrow it you must collect more \emph{real} data (longer
$N$) or trade resolution for a cleaner window (Chapter on windowing).
\end{intuitionbox}

\begin{trapbox}
\textbf{Two real, distinct uses of zero-padding, often confused:}
\begin{enumerate}
\item \emph{Spectral interpolation / nicer plots:} pad to $2N$, $4N$, $\dots$ to draw a
denser magnitude curve or to pin down a peak frequency between bins.
\item \emph{Linear convolution via the FFT:} to convolve two sequences correctly with
the FFT you must zero-pad \emph{both} to at least the sum of their lengths minus one,
otherwise the circular convolution the DFT computes will wrap around and corrupt the
result. (We return to this in the convolution chapter.)
\end{enumerate}
Reaching a power-of-two length for speed is a happy by-product of (1), not a third
purpose.
\end{trapbox}

\section{A short NumPy example}
\label{ch09:numpy}

Talk is cheap; here is a complete, correct script that synthesizes two tones, transforms
them, and reads amplitudes back using the calibration of \eqref{ch09:amp-cal}.

\begin{lstlisting}[language=Python, caption={Computing and reading a real FFT.}, label={ch09:lst-numpy}]
import numpy as np

fs = 128.0          # sampling rate, Hz
N  = 128            # number of samples (power of two -> fast)
t  = np.arange(N) / fs

# Signal: 1.0 amplitude at 12 Hz + 0.5 amplitude at 24 Hz + DC offset
x = 0.3 + 1.0*np.cos(2*np.pi*12*t) + 0.5*np.cos(2*np.pi*24*t)

X     = np.fft.rfft(x)                 # only non-redundant bins: 0..N/2
freqs = np.fft.rfftfreq(N, d=1/fs)     # frequency of each bin, in Hz

# Amplitude calibration: 2/N for interior bins, 1/N for DC and Nyquist
amp = (2.0/N) * np.abs(X)
amp[0] /= 2.0                          # DC bin has no mirror
if N % 2 == 0:
    amp[-1] /= 2.0                     # Nyquist bin has no mirror

# Report the strongest few bins
peaks = np.argsort(amp)[::-1][:3]
for k in sorted(peaks):
    print(f"{freqs[k]:6.1f} Hz   amplitude ~ {amp[k]:.3f}")

# Verify the FFT matches the brute-force DFT exactly (to rounding)
n = np.arange(N)
F = np.exp(-2j*np.pi*np.outer(n, n)/N)      # the DFT matrix
X_full = np.fft.fft(x)
assert np.allclose(X_full, F @ x)           # same numbers, ~N^2 vs N log N work
print("FFT matches direct DFT:", np.allclose(X_full, F @ x))
\end{lstlisting}

Running it prints frequencies near $0.0$, $12.0$, and $24.0$ Hz with amplitudes near
$0.3$, $1.0$, and $0.5$ -- the very signal we put in -- and confirms the FFT equals the
dense DFT matrix product. Two small but vital habits live in that listing: use
\texttt{rfftfreq} so you never hand-compute bin frequencies, and never trust a bin
magnitude as an amplitude until you have applied the $2/N$ scale.

\begin{labbox}
Modify the listing to time both methods on $N = 4096$:
\begin{lstlisting}[language=Python]
import time
x = np.random.randn(4096)
n = np.arange(4096); F = np.exp(-2j*np.pi*np.outer(n,n)/4096)
t0=time.perf_counter(); _=F@x;            t_dft=time.perf_counter()-t0
t0=time.perf_counter(); _=np.fft.fft(x);  t_fft=time.perf_counter()-t0
print(f"DFT {t_dft*1e3:.1f} ms   FFT {t_fft*1e3:.3f} ms   "
      f"speedup {t_dft/t_fft:.0f}x")
\end{lstlisting}
You should see the FFT beat the matrix multiply by two to three orders of magnitude --
and that is against NumPy's heavily optimized BLAS, which makes the matrix multiply far
faster than a textbook triple loop would be. Against naive C the gap is larger still.
\end{labbox}

\section{The inverse FFT comes for free}
\label{ch09:ifft}

The inverse DFT,
\begin{equation}
\label{ch09:idft}
x[n] = \frac{1}{N}\sum_{k=0}^{N-1} X[k]\,W_N^{-kn},
\end{equation}
differs from the forward transform only by a sign in the exponent and a $1/N$ scale. So
the \emph{same} butterfly machinery computes it -- conjugate the inputs, run a forward
FFT, conjugate the outputs, and divide by $N$; or simply use the library's
\texttt{ifft}. The practical consequence is large: any operation easy in the frequency
domain (filtering, convolution, correlation) becomes a round trip -- FFT, multiply,
inverse FFT -- each leg $O(N\log N)$. This ``transform, operate, transform back'' pattern
is the workhorse of the rest of the book.

\begin{applicationbox}
\textbf{Fast convolution.} Filtering a length-$N$ signal with a length-$M$ kernel costs
$O(NM)$ by direct convolution. Pad both to $N+M-1$, FFT each, multiply the spectra
pointwise, and inverse-FFT: total $O((N+M)\log(N+M))$. For long filters this is
dramatically faster, and it is how audio reverb plugins convolve a signal with a
multi-second impulse response in real time. The convolution theorem (later chapter) is
what licenses the pointwise multiply.
\end{applicationbox}

\section{Toward the wrist tracker}
\label{ch09:wrist}

\begin{goalbox}
Our gesture controller will carry an inertial measurement unit (IMU: accelerometer plus
gyroscope) on the forearm, and we want to react to motion in real time -- ideally with
latency low enough to feel instant while playing. A tremor, a tap, a wrist flick, the
rhythmic shake of a shaker gesture: these have signatures in the \emph{frequency} content
of the IMU streams, not just their raw values. The plan is to buffer a short window of
samples, FFT it, and read features off the spectrum (dominant frequency, band energies).

Make it concrete. Sample the IMU at $f_s = 200$ Hz and use a window of $N = 256$ samples.
Then:
\begin{itemize}
\item Window duration $= N/f_s = 256/200 = 1.28$ s of history per FFT.
\item Bin spacing $\Delta f = f_s/N = 200/256 \approx 0.78$ Hz -- fine enough to tell a
$3$ Hz tremor from a $5$ Hz one.
\item Useful band: $0$ to Nyquist $= f_s/2 = 100$ Hz, far above human limb motion (mostly
under $20$ Hz), so we have generous headroom and can even downsample.
\item Cost: $\tfrac{N}{2}\log_2 N = 128 \times 8 = 1024$ butterflies -- a few thousand
floating-point operations, trivial for an FFT a few hundred times per second even on a
Cortex-M microcontroller. The naive DFT at $N^2 = 65536$ would be far heavier and might
miss the real-time budget.
\end{itemize}
If $1.28$ s of latency is too sluggish (it is, for percussive triggers), shrink the
window or slide it with overlap; that is the short-time Fourier transform of the next
chapter. The FFT is what makes any of this affordable on a wrist-worn battery device.
\end{goalbox}

\begin{figure}[H]
\centering
\begin{tikzpicture}
\begin{axis}[
  width=12cm, height=4.6cm,
  xlabel={time (s)}, ylabel={accel.\ (a.u.)},
  domain=0:1.28, samples=180,
  xmin=0, xmax=1.28, ymin=-2.2, ymax=2.2,
  xtick={0,0.32,0.64,0.96,1.28},
  grid=both, grid style={gray!16},
  axis lines=box,
]
% a 4 Hz shaker gesture plus a small higher-frequency wobble
\addplot[sigblue, very thick] {1.4*sin(deg(2*pi*4*x)) + 0.35*sin(deg(2*pi*9*x))};
\end{axis}
\end{tikzpicture}
\caption{A simulated forearm accelerometer trace during a rhythmic shake: a strong $4$ Hz
oscillation with a weaker $9$ Hz wobble, over one $N=256$ window at $f_s=200$ Hz
($1.28$ s). The FFT of this block turns the tangle of wiggles into two clean peaks, from
which the controller reads the dominant gesture rate.}
\label{ch09:fig-imu}
\end{figure}

The frequency picture is what makes this tractable: Figure \ref{ch09:fig-imu} shows a raw
IMU trace whose rhythm is obvious to a spectrum but awkward to read off the time series.

\begin{goalbox}
\textbf{Hard limits -- read this twice.} An FFT of a wrist or forearm IMU tells you
about the \emph{gross motion of the whole limb}: is it shaking, at what rate, how
vigorously, and coarse posture changes. It \emph{cannot} recover individual finger joint
angles. The kinematic information about finger flexion simply is not present in the
arm-mounted accelerometer/gyro signal in any reliable form; no transform conjures
information the sensor never captured. To sense fingers you need a different modality:
surface electromyography (sEMG) or force myography (FMG) over the forearm muscles can
discriminate a handful of discrete finger/grip gestures (commercial armbands such as the
former Myo, and the research behind Meta's wrist sEMG work, demonstrate this), and
\emph{continuous} finger pose realistically needs sensors on the fingers themselves
(instrumented gloves or rings). So: FFT on the IMU is excellent for rhythm, taps, flicks,
and activity state -- treat finger-level control as a separate, harder problem requiring
EMG/FMG or finger-mounted sensing, and say so plainly in any demo.
\end{goalbox}

\begin{buildbox}
On an Arduino-class board, a fixed-point or float radix-2 FFT is cheap if you respect a
few rules.
\begin{itemize}
\item \textbf{Pick a power-of-two window} ($128$, $256$, $512$). Non-power-of-two sizes
drag in slower mixed-radix code or padding.
\item \textbf{Use a vendor library, not your own bit-reversal.} On ARM Cortex-M, ARM's
\texttt{CMSIS-DSP} provides \texttt{arm\_rfft\_fast\_f32} -- a tuned real FFT with the
twiddle tables baked in. On an ESP32, the ESP-DSP library has equivalents.
\item \textbf{Remove DC first.} Subtract the running mean of the accelerometer before the
FFT, or bin $0$ will dominate and swamp the motion bands you care about.
\item \textbf{Window the block} (Hann) before transforming to reduce spectral leakage
from the abrupt block edges; this matters more than zero-padding for clean peaks.
\item \textbf{Budget the memory.} A $256$-point float FFT needs roughly $256$ complex
words of scratch plus a twiddle table -- a few kilobytes, fine on most M-class parts.
\end{itemize}
A practical loop: fill a $256$-sample ring buffer at $200$ Hz, every $64$ new samples copy
it out, remove mean, apply Hann, \texttt{arm\_rfft\_fast\_f32}, take magnitude, then
classify from band energies. That overlap gives a fresh spectrum every $64/200 = 0.32$ s
while keeping the $0.78$ Hz resolution of a $256$-point window.
\end{buildbox}

\section{Practical notes and gotchas}
\label{ch09:practical}

A short field guide to the mistakes people actually make.

\begin{itemize}
\item \textbf{Spectral leakage.} A tone whose frequency falls \emph{between} bins does
not sit in one bin; its energy smears across neighbors (the sinc skirts of the window).
Apply a window function (Hann, Hamming, Blackman) before the FFT to suppress the side
lobes -- at the cost of slightly wider main lobes. This is a resolution-versus-leakage
trade, covered fully in the windowing chapter.
\item \textbf{Power spectrum scaling.} For energy/power, square the magnitude and apply
the right normalization for your window; do not compare raw $\abs{X[k]}^2$ across
different window functions or lengths without renormalizing.
\item \textbf{Real vs.\ complex transform.} For real sensor data use \texttt{rfft}; it is
about twice as fast and returns exactly the non-redundant half. Only use the full complex
\texttt{fft} when your data is genuinely complex (e.g.\ IQ radio samples).
\item \textbf{In-place and ordering.} If you call a low-level FFT, know whether it expects
or returns bit-reversed order, and whether it overwrites the input. High-level libraries
hide this; embedded ones often do not.
\item \textbf{Numerical accuracy.} The FFT is not only faster than the naive DFT, it is
usually \emph{more accurate}: $O(\log N)$ rounding steps accumulate per output instead of
$O(N)$ in the long naive sum. Speed and accuracy point the same way here.
\end{itemize}

\begin{recoverybox}
The whole chapter in five beats: (1) the DFT is an $N\times N$ matrix multiply, $O(N^2)$.
(2) The twiddle factor's half-turn symmetry means a size-$N$ DFT contains two size-$N/2$
DFTs. (3) Splitting input by even/odd and recombining with a butterfly turns that into
$O(N\log N)$. (4) The recursion's index shuffle is bit-reversal. (5) To \emph{use} an
FFT: bin $k$ sits at $kf_s/N$ Hz, the top half mirrors the bottom for real data, scale by
$2/N$ for amplitude, and zero-padding interpolates the plot without adding resolution.
\end{recoverybox}

\section{Exercises}
\label{ch09:exercises}

\begin{exercise}
\label{ch09:ex-count}
For $N = 4096$, compute the exact number of complex multiplications used by (a) the naive
DFT ($N^2$) and (b) the radix-2 FFT ($\tfrac{N}{2}\log_2 N$). What is the ratio? If one
complex multiply takes $5$ ns, estimate the wall-clock time of each.
\end{exercise}

\begin{exercise}
\label{ch09:ex-byhand}
Compute the $4$-point DFT of $x = (0,1,2,3)$ by hand using the radix-2 split of Section
\ref{ch09:split} (find $E$, $O$, then the butterflies). Verify $X[0]$ equals the sum of
the samples and check conjugate symmetry $X[3]=\overline{X[1]}$.
\end{exercise}

\begin{exercise}
\label{ch09:ex-bitrev}
Write out the bit-reversal permutation for $N = 16$ (a $4$-bit index). Which input index
ends up in slot $9$? Confirm your answer by reversing the bits of $9 = 1001_2$.
\end{exercise}

\begin{exercise}
\label{ch09:ex-symmetry}
Prove that if $x[n]$ is real then $X[N-k] = \overline{X[k]}$. Start from the DFT
definition \eqref{ch09:dft-defn} and use $W_N^{(N-k)n} = \overline{W_N^{kn}}$. What does
this say about how many of the $N$ outputs are independent?
\end{exercise}

\begin{exercise}
\label{ch09:ex-binfreq}
You sample at $f_s = 8$ kHz and take a $1024$-point FFT. (a) What frequency does bin $200$
represent? (b) What is the bin spacing $\Delta f$? (c) Which bin is closest to a $1$ kHz
tone, and what is the residual error in Hz?
\end{exercise}

\begin{exercise}
\label{ch09:ex-amp}
A real cosine of amplitude $3.0$ at exactly $50$ Hz is sampled at $f_s = 200$ Hz, $N =
200$. (a) Which bin does it land in? (b) What magnitude $\abs{X[k]}$ do you expect in that
bin before calibration? (c) Apply \eqref{ch09:amp-cal} to recover the amplitude.
\end{exercise}

\begin{exercise}
\label{ch09:ex-zeropad}
You have $N = 100$ samples of a signal sampled at $f_s = 1000$ Hz containing two tones at
$210$ Hz and $215$ Hz. (a) What is $\Delta f$? Are the tones resolved? (b) If you
zero-pad to $1000$ points, does the FFT now separate them into two distinct peaks?
Explain in terms of resolution versus interpolation. (c) What \emph{would} let you
resolve them?
\end{exercise}

\begin{exercise}
\label{ch09:ex-recurrence}
The FFT satisfies $T(N) = 2T(N/2) + cN$ for some constant $c$ (the $cN$ is the butterfly
combine step). Solve this recurrence by unrolling it $\log_2 N$ times and show the
solution is $T(N) = cN\log_2 N + T(1)\cdot N$. Identify which term dominates for large
$N$.
\end{exercise}

\begin{exercise}
\label{ch09:ex-numpy}
Adapt the listing in Section \ref{ch09:numpy} to a noisy signal: add
\texttt{0.2*np.random.randn(N)} to \texttt{x}. Does the peak detection still find the two
tones? Now reduce the tone amplitudes until a peak disappears into the noise floor. What
does this tell you about detecting weak gestures in a noisy IMU stream?
\end{exercise}

\begin{exercise}
\label{ch09:ex-wrist}
For the wrist tracker, you want at least $0.5$ Hz frequency resolution and you sample the
IMU at $f_s = 200$ Hz. (a) What is the minimum window length $N$ (round up to a power of
two)? (b) How many seconds of latency does that window impose? (c) If that latency is too
high for triggering a drum hit, name two changes (from this chapter and the next) that
reduce perceived latency, and state what each costs.
\end{exercise}

\begin{exercise}
\label{ch09:ex-limits}
In two or three sentences, explain to a teammate why running an FFT on the forearm IMU
\emph{cannot} give per-finger joint angles, and what sensing modality you would add if
the application genuinely needed finger-level control. Be specific about the limitation,
not vague.
\end{exercise}
